\documentclass[12pt]{article}
\usepackage{amsmath, amssymb}
\usepackage{fontspec}
\usepackage[pdfencoding=unicode,psdextra]{hyperref}
\hypersetup{pdfauthor={Hazime Sendo},pdftitle={DGUT 1.0: elementary structure and provisional dynamics},pdfsubject={Verification note for DGUT},colorlinks=false,hidelinks}

\setmainfont{Noto Serif CJK JP}
\usepackage[margin=2.5cm]{geometry}

\title{DGUT 1.0: 初等構造と暫定力学}
\author{仙道一 / Hazime Sendo}
\date{}

\begin{document}
\sloppy
\maketitle

\section*{第0章：初等構造（part0\_v10 準拠）}

\subsection*{0.1 ガウス平面と共役作用子}

本稿ではガウス平面 $\mathbb{C}$ を用い，


\[
Z = \phi + i\pi
\]


と書く。共役作用子 $\iota$ を


\[
\iota(Z) = \bar Z
\]


と定義する。これは involution であり，


\[
\iota\circ\iota = \mathrm{id}
\]


が成り立つ。

\subsection*{0.2 変化率 $A$ と符号分解（Prop.~sign）}

原稿の定義に従い，


\[
A(Z) = \langle \nabla F(\phi),\, \pi \rangle
\]


とする。ここから


\[
A(\iota Z) = -A(Z)
\]


が導かれる。

符号領域を


\[
\Sigma_+ = \{Z : A(Z) > 0\},\qquad
\Sigma_- = \{Z : A(Z) < 0\}
\]


と定義すると，$\iota$ により


\[
\iota(\Sigma_+) = \Sigma_-,\qquad
\iota(\Sigma_-) = \Sigma_+
\]


が成り立つ。

\subsection*{0.3 写像圧力 $\Lambda$（\textbf{最小 Lipschitz 定数}）}

写像 $f:\mathbb{C}\to\mathbb{C}$ に対し，
原稿の定義どおり最小 Lipschitz 定数を


\[
\Lambda(f)
= \sup_{x\neq y}\frac{|f(x)-f(y)|}{|x-y|}
\]


とする。

\paragraph{Remark.}
原点からの利得


\[
\sup_{Z\neq 0}\frac{|f(Z)|}{|Z|}
\]


とは一般に異なる。
たとえば $f(Z)=Z+1$ では $\Lambda(f)=1$ だが，
原点利得は $f(0)\neq 0$ のため無限大となる。

\subsection*{0.4 点ごとの上側 Lipschitz 定数 $\lambda_f$}

原稿の定義に従い，


\[
\lambda_f(x)
= \limsup_{y\to x}\frac{|f(y)-f(x)|}{|y-x|}
\]


とする。微分可能性は仮定しない。

一般には


\[
\sup_x \lambda_f(x) \le \Lambda(f)
\]


が成り立つが，
$\mathbb{C}$ は length space であるため，
原稿の命題により


\[
\sup_x \lambda_f(x) = \Lambda(f)
\]


となる。

\paragraph{Remark.}
$\mathbb{R}\setminus\{0\}$ は length space でないため，
$\sup_x \lambda_f(x) = \Lambda(f)$ が成り立つとは限らず，
実際に崩れる例が得られる。
この場合，$f=\mathrm{sign}$ では $\lambda_f(x)=0$ であっても，
$0$ をまたぐと $\Lambda(f)=\infty$ となる。

\subsection*{0.5 左入れ子再帰（dyad）と上界}

原稿の定義どおり，
2 引数写像（dyad）$f:\mathbb{C}^2\to\mathbb{C}$ の
左入れ子再帰を


\[
f^{[1]}(x_1,x_2)=f(x_1,x_2),\qquad
f^{[n]}(x_1,\dots,x_{n+1})
= f\bigl(f^{[n-1]}(x_1,\dots,x_n),\, x_{n+1}\bigr)
\]


と定義する。

$\mathbb{C}^{n+1}$ には和計量


\[
d(x,y) = \sum_{i=1}^{n+1} |x_i - y_i|
\]


を入れる。このとき原稿の命題として，


\[
\Lambda\bigl(f^{[n]}\bigr)
\le
\Lambda(f)\,\max(1,\Lambda(f))^{n-1}
\]


が成り立つ。

\paragraph{Remark.}
1 変数写像の反復 $f^n$ とは別の概念である。
反復では $\Lambda(f^n)\le \Lambda(f)^n$ が成り立つ。

\section*{Provisional Hypothesis P1（\textbf{暫定}）：最小線形力学}

part0 の Hypothesis 4.1 は運動方程式を述べない。
以下は part0 には含まれない新しい仮説（暫定）であり，part0 の Hypothesis 4.1 とは別のものとして P1 と呼ぶ。

\[
\dot Z = aZ + b\bar Z - \gamma Z.
\]

ここで，
\begin{itemize}
\item $a\in i\mathbb{R}$：保存的成分（\textbf{回転のみ}）
\item $b\in\mathbb{C}$：共役作用（鏡像）
\item $\gamma>0$：減衰率（実部を担う）
\end{itemize}

\paragraph{Remark 1（散逸にならない反例）.}
$a=a_2 i$ と書くと，固有値は
\[
(-\gamma)\pm\sqrt{|b|^2-a_2^2}
\]
である。ただし，これは $|b|\ge|a_2|$ のときの実根であり，
$|b|<|a_2|$ のときは複素共役の対となり，
実部は常に $-\gamma$ となる。
いずれにせよ，$\gamma>0$ でも成長する場合がある。
例：$a=0$, $b=2$, $\gamma=1$ では固有値 $+1$。

\paragraph{Remark 2（非正規性）.}
$M$ を
\[
M: Z \mapsto (a-\gamma)Z + b\bar Z
\]
とする $\mathbb{R}$-線形写像とする。
このときフローを
\[
\Phi_t = \mathrm{e}^{tM}
\]
と書き，$t=1$ の値を $\mathrm{e}^M$ と書く。
ここでのノルム $\|\mathrm{e}^{tM}\|$ は，
$\mathbb{C}\cong\mathbb{R}^2$ のユークリッドノルムに対する
作用素ノルム（最大特異値）である\cite{HJ}。

固有値の実部が負でも，
\[
\|\mathrm{e}^{M}\|>1
\]
となる場合がある。
例：$a=i$, $b=0.3$, $\gamma=0.1$ では
実部 $-0.1$ だが $\|\mathrm{e}^M\|=1.166$。

線形フローであるため，
\[
\Lambda(\Phi_t) = \|\mathrm{e}^{tM}\|
\]
が成り立つ。これは原稿の Prop.~6.1(c)
「$\Lambda_D(\varphi_t)<1$ なら $\Pi_Q$ は定数」
に接続するが（ここでの $\Phi_t$ を Prop.~6.1 の $\varphi_t$ とみなすことは仮定である），
$\Pi_Q$ 自体は未定義のままである。

\paragraph{Remark 3（$\iota$ 対称性と時間反転）.}
$\iota\circ f = f\circ\iota$ は，$\bar a=a$ かつ $\bar b=b$ と同値である。
$a\in i\mathbb{R}$ の下では $a=0$ かつ $b\in\mathbb{R}$ となり，回転は残らない。

時間反転型の対称性
\[
\iota\circ f = -f\circ\iota
\]
は，$a-\gamma$ が純虚数で $b$ も純虚数のときに限って成り立つ。
$a\in i\mathbb{R}$ の下では，これは $\gamma=0$ かつ $b\in i\mathbb{R}$ と同値である。
したがって $\gamma>0$ では，許される全てのパラメータで
この時間反転対称性は成り立たない。

\paragraph{Remark 4（符号分解との整合）.}
本 ODE は $F$ と $A=\langle\nabla F,\pi\rangle$ を含まないため，
符号分解との整合は未確立である。

\section*{初等構造の数値例}

本節の数値例は，同梱のコード \texttt{src/dgut\_elementary.py}
（0.1--0.5 の初等命題の例示）および
\texttt{src/dgut\_ode\_check.py}
（Provisional Hypothesis P1 の暫定 ODE の線形代数的性質：Remark 1--3 と
$\Lambda(\Phi_t)=\|\mathrm{e}^{tM}\|$）による。
これらは初等命題と暫定 ODE の性質の数値例であり，
「確認された」ほどの強さは持たない。
$\Pi_Q$，状態割当 $\chi$，相の対応は原稿のプレースホルダーであり，
検証の対象外である。


\begin{thebibliography}{9}
\bibitem{HJ}
R.~A.~Horn, C.~R.~Johnson,
\textit{Matrix Analysis},
Cambridge University Press (1985).
\end{thebibliography}

\end{document}
